\chapter{Backend Implementation}

\section{Technology Stack}

The backend of Campus-Sync is built using a modern, scalable technology stack designed to handle the complex requirements of an educational management platform:

\begin{itemize}
    \item \textbf{Express.js}: A minimal and flexible Node.js web application framework
    \item \textbf{MongoDB}: A document-based NoSQL database for flexible data storage
    \item \textbf{Mongoose}: An Object Data Modeling (ODM) library for MongoDB and Node.js
    \item \textbf{JWT (JSON Web Tokens)}: A standard for securely transmitting information between parties
    \item \textbf{Swagger/OpenAPI}: API documentation framework
    \item \textbf{Joi/Zod}: Validation libraries for data validation
    \item \textbf{Winston}: A logger for Node.js applications
    \item \textbf{Jest/Supertest}: Testing frameworks for unit and integration testing
    \item \textbf{Docker}: Containerization platform for consistent deployment
\end{itemize}

\section{Database Design}

The database schema is designed to support all aspects of the educational management system with proper relationships and normalization.

\subsection{Core Collections}

\subsubsection{Users Collection}

\begin{lstlisting}[language=JavaScript, caption=User Schema]
{
  _id: ObjectId,
  fullName: String,
  email: String,
  password: String,
  role: Enum['student', 'teacher', 'admin'],
  profilePicture: String,
  dateOfBirth: Date,
  phoneNumber: String,
  address: {
    street: String,
    city: String,
    state: String,
    zipCode: String,
    country: String
  },
  isActive: Boolean,
  lastLogin: Date,
  createdAt: Date,
  updatedAt: Date
}
\end{lstlisting}

\subsubsection{Students Collection}

\begin{lstlisting}[language=JavaScript, caption=Student Schema]
{
  _id: ObjectId,
  userId: ObjectId, // Reference to Users collection
  studentId: String,
  enrollmentDate: Date,
  branch: String,
  semester: Number,
  courses: [ObjectId], // References to Courses
  guardianInfo: {
    name: String,
    relationship: String,
    phoneNumber: String,
    email: String
  },
  academicRecords: [{
    courseId: ObjectId,
    semester: Number,
    grade: String,
    credits: Number
  }],
  createdAt: Date,
  updatedAt: Date
}
\end{lstlisting}

\subsubsection{Teachers Collection}

\begin{lstlisting}[language=JavaScript, caption=Teacher Schema]
{
  _id: ObjectId,
  userId: ObjectId, // Reference to Users collection
  employeeId: String,
  department: String,
  specialization: String,
  hireDate: Date,
  salary: Number,
  courses: [ObjectId], // References to Courses
  createdAt: Date,
  updatedAt: Date
}
\end{lstlisting}

\subsubsection{Courses Collection}

\begin{lstlisting}[language=JavaScript, caption=Course Schema]
{
  _id: ObjectId,
  code: String,
  name: String,
  description: String,
  credits: Number,
  semester: Number,
  branch: String,
  prerequisites: [String],
  instructor: ObjectId, // Reference to Teachers
  schedule: {
    days: [String],
    startTime: String,
    endTime: String,
    room: String
  },
  createdAt: Date,
  updatedAt: Date
}
\end{lstlisting}

\subsubsection{Assignments Collection}

\begin{lstlisting}[language=JavaScript, caption=Assignment Schema]
{
  _id: ObjectId,
  title: String,
  description: String,
  courseId: ObjectId, // Reference to Courses
  assignedBy: ObjectId, // Reference to Teachers
  dueDate: Date,
  maxPoints: Number,
  attachments: [String],
  submissions: [{
    studentId: ObjectId, // Reference to Students
    submittedAt: Date,
    fileUrl: String,
    points: Number,
    feedback: String
  }],
  createdAt: Date,
  updatedAt: Date
}
\end{lstlisting}

\subsection{Specialized Collections}

\subsubsection{Tasks Collection}

\begin{lstlisting}[language=JavaScript, caption=Task Schema]
{
  _id: ObjectId,
  userId: ObjectId, // Reference to Users
  title: String,
  description: String,
  dueDate: Date,
  priority: Enum['low', 'medium', 'high'],
  status: Enum['todo', 'in-progress', 'completed'],
  tags: [String],
  attachments: [String],
  createdAt: Date,
  updatedAt: Date
}
\end{lstlisting}

\subsubsection{Events Collection}

\begin{lstlisting}[language=JavaScript, caption=Event Schema]
{
  _id: ObjectId,
  title: String,
  description: String,
  startDate: Date,
  endDate: Date,
  location: String,
  organizer: ObjectId, // Reference to Users
  attendees: [ObjectId], // References to Users
  attachments: [String],
  isPublic: Boolean,
  createdAt: Date,
  updatedAt: Date
}
\end{lstlisting}

\subsubsection{Blog Collection}

\begin{lstlisting}[language=JavaScript, caption=Blog Schema]
{
  _id: ObjectId,
  title: String,
  content: String,
  author: ObjectId, // Reference to Users
  tags: [String],
  category: String,
  featuredImage: String,
  likes: [ObjectId], // References to Users
  comments: [{
    userId: ObjectId, // Reference to Users
    content: String,
    createdAt: Date
  }],
  published: Boolean,
  publishedAt: Date,
  createdAt: Date,
  updatedAt: Date
}
\end{lstlisting}

\subsubsection{Music Collection}

\begin{lstlisting}[language=JavaScript, caption=Music Schema]
{
  _id: ObjectId,
  title: String,
  artist: String,
  album: String,
  genre: String,
  duration: Number, // in seconds
  fileUrl: String,
  coverArt: String,
  uploadedBy: ObjectId, // Reference to Users
  plays: Number,
  likes: [ObjectId], // References to Users
  createdAt: Date,
  updatedAt: Date
}
\end{lstlisting}

\subsubsection{AI Assistant Collection}

\begin{lstlisting}[language=JavaScript, caption=AI Assistant Schema]
{
  _id: ObjectId,
  userId: ObjectId, // Reference to Users
  prompt: String,
  response: String,
  category: String, // 'academic', 'financial', 'general'
  createdAt: Date
}
\end{lstlisting}

\section{API Endpoints}

The backend provides a comprehensive RESTful API with endpoints organized by functionality:

\subsection{Authentication Routes}

\begin{verbatim}
POST /api/auth/register - User registration
POST /api/auth/login - User login
POST /api/auth/refresh - Refresh access token
POST /api/auth/logout - User logout
POST /api/auth/forgot-password - Password reset request
POST /api/auth/reset-password - Password reset
\end{verbatim}

\subsection{User Management Routes}

\begin{verbatim}
GET /api/users/profile - Get user profile
PUT /api/users/profile - Update user profile
PUT /api/users/password - Change password
DELETE /api/users/profile - Delete user account
\end{verbatim}

\subsection{Student Routes}

\begin{verbatim}
GET /api/students/dashboard - Student dashboard data
GET /api/students/timetable - Student timetable
GET /api/students/courses - Student courses
GET /api/students/grades - Student grades
GET /api/students/attendance - Student attendance
GET /api/students/progress - Student academic progress
\end{verbatim}

\subsection{Teacher Routes}

\begin{verbatim}
GET /api/teachers/dashboard - Teacher dashboard data
GET /api/teachers/timetable - Teacher timetable
GET /api/teachers/courses - Teacher courses
GET /api/teachers/students - Students in teacher's courses
POST /api/teachers/assignments - Create assignment
PUT /api/teachers/assignments/:id - Update assignment
POST /api/teachers/grades - Submit grades
\end{verbatim}

\subsection{Admin Routes}

\begin{verbatim}
GET /api/admin/dashboard - Admin dashboard data
GET /api/admin/students - Get all students
POST /api/admin/students - Create student
PUT /api/admin/students/:id - Update student
DELETE /api/admin/students/:id - Delete student
GET /api/admin/teachers - Get all teachers
POST /api/admin/teachers - Create teacher
PUT /api/admin/teachers/:id - Update teacher
DELETE /api/admin/teachers/:id - Delete teacher
GET /api/admin/courses - Get all courses
POST /api/admin/courses - Create course
PUT /api/admin/courses/:id - Update course
DELETE /api/admin/courses/:id - Delete course
\end{verbatim}

\section{Middleware and Utilities}

\subsection{Custom Middleware}

\begin{itemize}
    \item Authentication middleware to verify JWT
    \item Authorization middleware to check roles
    \item Validation middleware for request data
    \item Error handling middleware for consistent error responses
    \item Rate limiting middleware to prevent abuse
    \item Logging middleware for request tracking
\end{itemize}

\subsection{Utility Functions}

\begin{itemize}
    \item Password hashing and verification using bcrypt
    \item JWT token generation and verification
    \item Email sending utility
    \item File upload utility
    \item Date formatting utilities
    \item Pagination utility
\end{itemize}

\section{Security Considerations}

The backend implements comprehensive security measures to protect user data and system integrity:

\subsection{Authentication Security}

\begin{itemize}
    \item JWT-based authentication with secure token storage
    \item Password hashing with bcrypt
    \item Session management with refresh tokens
    \item Role-based access control (RBAC)
\end{itemize}

\subsection{Data Protection}

\begin{itemize}
    \item Input validation and sanitization
    \item Helmet.js for HTTP headers security
    \item CORS configuration
    \item Rate limiting to prevent abuse
    \item XSS protection
    \item CSRF protection
    \item Secure password storage
\end{itemize}

\subsection{API Security}

\begin{itemize}
    \item Environment-based configuration
    \item Secure API key management
    \item Database access control
    \item Regular security audits
\end{itemize}

\section{Performance Optimization}

Several optimization techniques are implemented to ensure fast response times and efficient resource usage:

\subsection{Database Optimization}

\begin{itemize}
    \item Proper indexing strategies for frequently queried fields
    \item Query optimization to minimize database load
    \item Connection pooling for efficient database connections
    \item Database sharding for scalability
\end{itemize}

\subsection{Caching}

\begin{itemize}
    \item Redis for session storage
    \item API response caching for frequently accessed data
    \item Database query caching
\end{itemize}

\subsection{Load Handling}

\begin{itemize}
    \item Clustering for multi-core systems
    \item Load balancing configuration
    \item Compression (gzip) for response data
    \item CDN for static assets
\end{itemize}

\section{Testing Strategy}

A comprehensive testing approach ensures the reliability and stability of the backend system:

\subsection{Test Types}

\begin{itemize}
    \item Unit tests for services and utilities
    \item Integration tests for API endpoints
    \item End-to-end tests for critical workflows
    \item Performance tests for load handling
\end{itemize}

\subsection{Testing Tools}

\begin{itemize}
    \item Jest for unit testing
    \item Supertest for API testing
    \item MongoDB in-memory server for database tests
    \item CI/CD pipeline integration
\end{itemize}

\section{Deployment and Monitoring}

The backend is designed for scalable deployment with comprehensive monitoring capabilities:

\subsection{Deployment}

\begin{itemize}
    \item Docker containerization for consistent environments
    \item Docker Compose for multi-container setup
    \item Environment-specific configurations
    \item CI/CD pipeline setup
\end{itemize}

\subsection{Monitoring}

\begin{itemize}
    \item Application performance monitoring (APM)
    \item Database performance monitoring
    \item Error tracking and logging
    \item Health checks and uptime monitoring
\end{itemize}

\subsection{Scaling}

\begin{itemize}
    \item Horizontal scaling strategies
    \item Load balancer configuration
    \item Database replication
    \item Caching strategies
\end{itemize}

This backend implementation provides a robust, secure, and scalable foundation for the Campus-Sync platform, ensuring reliable data management and API services for all user roles.